\chapter{Fontes, medidas e proveniência}\label{cap:fontes}
\section{Catálogo acessível da fonte Spotify}
A fonte inicial é o conjunto público identificado como \emph{Spotify\_1Million\_Tracks}, disponibilizado na plataforma Kaggle pela conta \texttt{amitanshjoshi} \citep{spotifydataset}. O arquivo local anterior a qualquer matching contém 1.159.764 linhas e o mesmo número de identificadores de faixa distintos. Sua auditoria verificou ausência de IDs faltantes ou repetidos. As observações têm gêneros registrados em 82 categorias e o campo \texttt{year} assume valores entre 2000 e 2023.

Esses registros definem o catálogo acessível deste estudo. Não se conhece um desenho probabilístico que tenha selecionado esse catálogo a partir de todo o Spotify. Sua dimensão e variedade não justificam tratar a fonte como censo da plataforma. A palavra catálogo, neste texto, designa esse arquivo congelado, identificado por caminho, tamanho e hash, e não a totalidade de músicas existentes no serviço.

A variável resposta é \texttt{spotify\_popularity}, derivada explicitamente do campo \texttt{popularity} da fonte. A documentação do Spotify descreve um escore de zero a cem, calculado principalmente com informação de execuções e recência, e informa que versões duplicadas podem ser avaliadas separadamente \citep{spotifydocs}. Essa definição não revela a data de medição de cada linha no arquivo utilizado. Não se trata de contagem de execuções nem de medida de popularidade no lançamento.

Há uma distinção entre datas documentais. O manifesto local registra reuso ou verificação do arquivo em 18 de agosto de 2026; isso não comprova a data de download original. Uma atualização upstream em 2023 consta da proveniência histórica, mas não equivale à data de coleta da resposta. A versão exata do conjunto Kaggle e o significado de \texttt{year} como ano original de lançamento permanecem não confirmados. A análise usa o campo como metadado fornecido, sem reconstruir datas por títulos ou letras.

Além da popularidade, os campos acústicos e musicais incluem danceability, energy, loudness, speechiness, acousticness, instrumentalness, liveness, valence, tempo, duração, tonalidade, modo e assinatura de tempo. São medidas da fonte, não extraídas novamente do áudio pelo projeto. Nomes de campos não substituem interpretação de suas escalas; tampouco sua presença autoriza atribuição causal.

\section{LFM-2b e descritores de letras: papel histórico}
O LFM-2b reúne eventos de escuta associados ao Last.fm. A derivação \emph{LFM2b Lyrics Descriptor Analyses} é outro objeto: disponibiliza descritores textuais e informações de músicas, com 582.759 registros completos segundo seus autores. A derivação usa letras identificadas como inglesas e distingue contagens de escuta no Last.fm de visualizações de letras no Genius \citep{parada2024data}.

No projeto, essa derivação participou de uma integração inicial com a fonte Spotify e motivou estudos de descritores. A integração histórica aceitou 111.758 pares, universo diferente do domínio recente. Os seus resultados, vetores e transformações não foram transportados automaticamente para a análise atual. Um \texttt{track\_id} comum não demonstra que a letra, a versão da gravação e o processamento textual sejam os mesmos.

Três variáveis precisam permanecer separadas: popularidade Spotify, visualizações de letras no Genius e eventos de escuta no Last.fm. Somente a primeira é resposta desta monografia. A informação LFM2b não é requisito de elegibilidade da amostra recente e seus descritores não preenchem silenciosamente as faixas que não integraram o matching antigo. Esse papel histórico é importante para a trajetória do estudo, mas não constitui um quinto bloco preditivo atual.

\section{LRCLIB e o snapshot textual}
As letras da versão atual foram recuperadas de um dump local do LRCLIB \citep{lrclib}. O arquivo identificado pelo projeto corresponde ao snapshot de 12 de outubro de 2026. Ele armazena metadados e representações textuais que podem conter letra simples, letra sincronizada, sinalização instrumental ou ausência de conteúdo utilizável. A existência de um registro não garante correspondência correta com uma faixa Spotify.

O armazenamento local permitiu consultas reproduzíveis sem enviar letras a serviços externos de inferência. Índices auxiliares por artista apoiaram a recuperação; suas contagens são contagens de registros indexados, não de letras aceitas na amostra. Nenhum índice ou varredura foi reconstruído para produzir esta monografia.

Os procedimentos de preparação do dump registraram tamanho, verificações de transporte e consistência SQLite. Esses controles detectam corrupção técnica, mas não são assinatura independente do publicador nem avaliação da qualidade semântica dos textos. A disponibilidade pública de uma base tampouco implica licença irrestrita sobre todas as letras. O pacote da monografia contém agregados, fontes de figuras e definições; não distribui letras, bases individuais, credenciais ou pesos de modelos.

\section{Unidades de observação e ligação}
A unidade inicial e de avaliação principal é o identificador de faixa Spotify. Um identificador não é uma composição universal: gravações, álbuns e versões podem gerar registros distintos. Uma composição operacional é uma chave de agrupamento construída no projeto; uma letra aceita é um conteúdo ligado por referência e hash; um componente operacional pode reunir várias faixas e relações. Essas unidades são usadas para finalidades diferentes e não devem ter suas contagens trocadas.

Os hashes documentam identidade de arquivos ou entradas normalizadas. Eles permitem detectar alteração ou reuso incompatível, mas não provam que uma letra pertença à gravação correta. A cadeia de proveniência informa de onde veio o conteúdo e qual regra o aceitou; a validade substantiva do matching continua exigindo evidência adicional.
